\ncertchapter{5}{Applications of Radioactivity}\chlabel{ch:nucapps}

A nucleus decays at a moment no one can predict, yet a population of nuclei decays on a schedule that
nothing can disturb: not heat, not pressure, not chemistry. That combination, a law that governs only the
statistics but governs them exactly, is what makes radioactivity useful. Everything needed is now in
place. \chref{ch:decay} gave the energetics of each decay and the particles it releases;
\chref{ch:decaylaw} gave the exponential law $N(t)=N_0e^{-\lambda t}$, where $N(t)$ is the number of
undecayed nuclei at time $t$, $N_0$ the number at $t=0$ and $\lambda$ the decay constant (probability of
decay per unit time, in \si{s^{-1}}), together with the half-life $\Thalf=\ln2/\lambda$, the mean life
$\tau=1/\lambda$ and the activity $\Act=\lambda N$ (decays per second, in becquerels);
\chref{ch:reactions} showed how nuclei are made and transformed in reactions.

This chapter puts these tools to work in four areas. Radioactive \emph{clocks} read the decay law
backwards to date bones, rocks and the solar system. \emph{Medicine} uses radioactive atoms as tracers
that the body cannot tell from ordinary ones, and as sources that kill cancer cells. Because the same
radiation that heals can harm, the chapter then sets up the measurement of radiation \emph{dose} and the
three rules of protection. Finally it closes a loop opened in \chref{ch:decaylaw}: a computer that knows
nothing except ``each nucleus has a fixed chance of decaying in each small time step'' reproduces the
exponential law unaided, and shows why the law is smooth only for large numbers.

The limits are stated as they arise. Every date rests on assumptions about the sample's history; every
dose figure for low exposures rests on an extrapolation; every shielding formula assumes a narrow beam.
The chapter, and the volume, end with a section that combines ideas from all five chapters, and with a
look back over the whole volume.

\section{Radioactive Clocks}\label{sec:nucapp-clocks}

A clock needs a process whose rate is known and cannot be tampered with. Chemical and biological rates
change with temperature; the decay constant does not, because the energies that hold a nucleus together
(millions of electronvolts) dwarf anything that heat or chemistry can supply (a few electronvolts). A
radioactive isotope is therefore an ideal clock. To read it, two facts are needed: how much of the
parent is left now, and either how much there was at the start or how much daughter has been made.

\begin{derivbox}{The dating equation}
Call the parent X and the daughter Y. Suppose a sample formed at $t=0$ with $N_{X,0}$ atoms of X and no
atoms of Y, and that it has stayed \emph{closed}: no atoms of either kind have entered or left. After a
time $t$ the number of parent atoms left is
\[
N_X=N_{X,0}\,e^{-\lambda t},
\]
with $\lambda$ the decay constant of X. Every parent that decayed is now a daughter, provided the
daughter is stable (or the intermediate members of a chain are short-lived), so $N_Y=N_{X,0}-N_X$. The
starting number $N_{X,0}$ is not known, but it can be eliminated: $N_{X,0}=N_X+N_Y$. Then
\[
N_X=(N_X+N_Y)\,e^{-\lambda t}\quad\Longrightarrow\quad e^{\lambda t}=1+\frac{N_Y}{N_X}.
\]
Taking logarithms (\chref{ch:toolkit}, Section~\ref{sec:explog}) and using $\lambda=\ln2/\Thalf$,
\[
t=\frac1\lambda\ln\!\Big(1+\frac{N_Y}{N_X}\Big)=\Thalf\log_2\!\Big(1+\frac{N_Y}{N_X}\Big).
\]
Check the limits: with no daughter, $N_Y=0$ and $t=0$; with equal numbers, $N_Y/N_X=1$ and
$t=\Thalf$, as it must be, since half the parent has then become daughter.
\end{derivbox}

\begin{keybox}{Radioactive dating}
For a closed sample that began with no daughter,
\[
t=\Thalf\log_2\!\Big(1+\frac{N_Y}{N_X}\Big).
\]
If the starting amount of parent is known by some other argument (as in radiocarbon dating), compare
the present amount or activity with it directly:
$t=\Thalf\log_2(N_{X,0}/N_X)=\Thalf\log_2(\Act_0/\Act)$.
\end{keybox}

Only a \emph{ratio} of two present-day numbers is needed, and a mass spectrometer measures ratios of
atoms precisely. Figure~\ref{fig:nucapp-dating} shows why the method reaches so far back. The ratio
$N_Y/N_X=2^{t/\Thalf}-1$ is $1,3,7,15,\dots$ after $1,2,3,4,\dots$ half-lives: it grows without limit,
so even a sample many half-lives old has a large, easily measured ratio. The practical limit is set
instead by the parent: once too little of it survives, its amount cannot be measured.

\begin{figure}[htbp]
\centering
\tikzfig{nucapp_dating}
\figcaption{fig:nucapp-dating}{A closed sample with no daughter at the start. The parent fraction
halves each half-life while the daughter fraction rises to fill the gap, so the two always add to $1$.
Their ratio $N_Y/N_X$ (dashed, right axis) doubles-plus-one each half-life ($1$, $3$, $7$, $15$) and
never levels off; this is what lets one measured ratio give an age.}
\end{figure}

\subsection{Radiocarbon}\label{sec:nucapp-c14}

Organic remains (wood, bone, charcoal, cloth) are the commonest things an archaeologist wants to date,
and they contain carbon. Carbon-14 is the clock. High in the atmosphere, neutrons released by cosmic rays
convert nitrogen into carbon-14 by the reaction
\[
n+\iso{14}{N}\to\iso{14}{C}+p,
\]
a reaction of the kind treated in \chref{ch:reactions}. The carbon-14 forms carbon dioxide and mixes
through the air. It decays by $\beta^-$ emission back to nitrogen,
$\iso{14}{C}\to\iso{14}{N}+e^-+\bar\nu_e$, with $\Thalf=\SI{5730}{\year}$. Production at a steady rate
balanced by decay is exactly the situation of constant production in \chref{ch:decaylaw}: the amount
settles at the level where the decay rate equals the production rate. The result is that about $1.2$
atoms in every $10^{12}$ atoms of atmospheric carbon are carbon-14. Plants take in this carbon dioxide by
photosynthesis and animals eat the plants, so every living thing shares the same carbon-14 fraction and
therefore the same activity per gram of carbon (the \emph{specific activity}).

\begin{derivbox}{The specific activity of living carbon}
One gram of carbon (molar mass \SI{12.01}{g.mol^{-1}}) contains $N_A/12.01=6.022\times10^{23}/12.01
=5.01\times10^{22}$ atoms, of which a fraction $1.2\times10^{-12}$ is carbon-14:
$N=6.0\times10^{10}$ atoms. The decay constant, with $\SI{1}{\year}=\SI{3.156e7}{s}$, is
\[
\lambda=\frac{\ln2}{\Thalf}=\frac{0.693}{5730\times3.156\times10^{7}\ \si{s}}=\SI{3.83e-12}{s^{-1}}.
\]
So the activity of one gram of living carbon is
\[
\Act=\lambda N=3.83\times10^{-12}\times6.0\times10^{10}=\SI{0.23}{\becquerel},
\]
that is, $0.23\times60\approx13.8$ decays per minute per gram of carbon. A tiny activity: a gram of
living carbon gives fewer than one decay every four seconds, which is why radiocarbon laboratories now
count carbon-14 atoms directly with an accelerator mass spectrometer rather than wait for them to decay.
\end{derivbox}

At death, intake stops. No new carbon-14 arrives and what is present decays away, while the stable
carbon-12 stays. Comparing a sample's specific activity with the living value gives its age through
$t=\Thalf\log_2(\Act_\text{living}/\Act_\text{sample})$, the second form of the dating result. The whole
cycle is summarised in Figure~\ref{fig:nucapp-c14cycle}.

\begin{figure}[htbp]
\centering
\tikzfig{nucapp_c14cycle}
\figcaption{fig:nucapp-c14cycle}{The radiocarbon cycle. Read the top row as a steady state: cosmic rays
keep topping up the atmosphere's carbon-14 while it decays, and living things share it. The arrow
marked ``death'' starts the clock (heavy boxes): from then on carbon-14 only decays, with
$\Thalf=\SI{5730}{\year}$, and the present activity measures the time elapsed.}
\end{figure}

\begin{exbox}{An ancient hearth}\label{ex:nucapp-hearth}
Charcoal from an ancient hearth shows $3.45$ decays per minute per gram of carbon; living wood shows
$13.8$. How old is the hearth?

\stepn{1} \textbf{Compare on the same basis.} Both activities are per gram of carbon, so their ratio is
the ratio of carbon-14 atoms: $\Act_\text{living}/\Act_\text{sample}=13.8/3.45=4=2^2$.

\stepn{2} \textbf{Count half-lives.} A factor of $4$ is two halvings, so $t=2\Thalf=2\times5730
=\SI{11460}{\year}$.
\ans{$\SI{1.15e4}{\year}$}
\end{exbox}

The method has limits, and each is instructive. Beyond about \SI{50000}{\year} (nearly nine half-lives,
leaving $2^{-8.7}\approx0.2\%$ of the original carbon-14) too little remains to measure against the
background. The production rate has not been perfectly steady, because the cosmic-ray flux varies with the
Sun's activity and the Earth's magnetic field; dates are corrected with a calibration curve built from tree
rings counted year by year. And human activity has disturbed the balance in both directions: burning
fossil fuel, whose carbon-14 decayed away long ago, dilutes the atmospheric fraction, while the atmospheric
nuclear-weapons tests of the 1950s and early 1960s nearly doubled it for a time. Samples are therefore
compared with the living standard as it was before 1950.

\subsection{Rocks and the Age of the Earth}\label{sec:nucapp-rocks}

Carbon-14 is far too short-lived for geological time. Rocks are dated with long-lived parents, using the
parent--daughter form of the dating result: uranium-238 decays through a chain to lead-206
($\Thalf=\SI{4.468e9}{\year}$), potassium-40 partly to argon-40 ($\Thalf=\SI{1.248e9}{\year}$) and
rubidium-87 to strontium-87 ($\Thalf=\SI{4.97e10}{\year}$). The uranium chain has fourteen steps, but
every intermediate member is short-lived compared with uranium-238, so on geological time scales each
uranium atom that decays becomes a lead atom almost at once and the derivation applies unchanged.

The hardest assumption to secure is ``no daughter at the start''. Nature supplies a mineral that secures
it: crystals of zircon accept uranium readily when they grow but exclude lead, so any lead found in a zircon
was made there by decay (it is \emph{radiogenic}). The oldest zircons, from the Jack Hills of Western
Australia, are $4.4$ billion years old. Meteorites, leftovers of the material from which the planets
formed, give the age of the solar system, $4.57$ billion years.

\begin{exbox}{Uranium--lead}\label{ex:nucapp-zircon}
A zircon contains $0.50$ atoms of lead-206 for every atom of uranium-238. How old is it?

\stepn{1} \textbf{Identify the form.} A daughter-to-parent ratio in a closed crystal that began with no
lead: $t=\Thalf\log_2(1+N_Y/N_X)$.

\stepn{2} \textbf{Compute.} $t=\SI{4.468e9}{\year}\times\log_2(1.50)=4.468\times10^{9}\times0.585
=\SI{2.61e9}{\year}$, a little over half a half-life, consistent with a ratio below $1$.
\ans{$\SI{2.61e9}{\year}$}
\end{exbox}

Potassium-40 adds one complication, which is \emph{branching} (\chref{ch:decaylaw}): it has two decay
channels. In $10.72\%$ of its decays it becomes argon-40 (by electron capture or $\beta^+$ emission); in
the rest it becomes calcium-40 by $\beta^-$ decay, and calcium is so common in rocks that the radiogenic part
cannot be picked out. Only argon is useful. The parent still disappears at the rate set by its \emph{total}
decay constant $\lambda$, so the number of parents that have decayed is $N_{X,0}-N_X=N_X(e^{\lambda t}-1)$,
and only the fraction $f=0.1072$ of them became argon:
\[
\frac{N_Y}{N_X}=f\left(e^{\lambda t}-1\right)\quad\Longrightarrow\quad
t=\Thalf\log_2\!\Big(1+\frac{N_Y}{fN_X}\Big).
\]
Argon is a gas and escapes from molten rock, so the clock starts when a lava flow solidifies, which is why
the potassium--argon method dates volcanic events.

\begin{trapbox}{the assumptions behind every date}
The dating equation assumes a \emph{closed system} (no parent or daughter gained or lost since the clock
started) and a known starting amount of daughter (often zero). A rock that was later reheated and lost some
of its argon will look too young; one that absorbed uranium-bearing water will look wrong in a way that is
hard to predict. Geologists guard against this by checking that different parent--daughter pairs in the
same rock give the same age. A single date from a single method is a measurement with hidden assumptions,
not a fact.
\end{trapbox}

\section{Radioactivity in Medicine}\label{sec:nucapp-medicine}

Dating used radioactivity as a clock. Medicine uses two other properties: a radioactive atom can be
\emph{followed} from outside the body, because its radiation escapes, and its radiation can \emph{kill}
cells, because it ionises. Diagnosis uses the first property, therapy the second.

\subsection{Diagnosis with Tracers}

A radioactive isotope has the same electrons, and so the same chemistry, as a stable isotope of the same
element. The body cannot tell them apart and sends the radioactive atom wherever it would send the stable
one; the gamma rays that escape reveal where it went. This idea, the \emph{radioactive tracer}, needs an
isotope that emits gamma rays (which leave the body, unlike $\alpha$ and $\beta$ particles), with an energy
high enough to escape tissue but low enough to be stopped in a detector, and a half-life of hours, long
enough to do the scan but short enough not to linger. Technetium-99m fits all three. It is an isomer,
a long-lived excited state of technetium-99 (\chref{ch:decay}), which decays by emitting a single
\SI{140}{keV} gamma ray with $\Thalf=\SI{6.0}{h}$. Attached to different carrier molecules, it images
bones, the heart, the kidneys and the brain.

Short half-lives create a logistical problem: the isotope decays on the way to the patient, so the
supplier must send more than the patient needs. Working backwards through the decay law, the activity
needed at dispatch, a time $t$ before use, is
\[
\Act_\text{dispatch}=\Act_\text{use}\times2^{t/\Thalf}.
\]

\begin{exbox}{Ordering technetium-99m}\label{ex:nucapp-tc99m}
A hospital needs \SI{500}{\mega\becquerel} of technetium-99m ($\Thalf=\SI{6.0}{h}$) at 9:00\,a.m. The
dose is prepared at 9:00\,p.m.\ the previous evening. What activity must be prepared?

\stepn{1} \textbf{Elapsed time in half-lives.} \SI{12}{h} is $12/6.0=2$ half-lives.

\stepn{2} \textbf{Work backwards.} The activity falls by $2^2=4$ in that time, so the dose must start at
$500\times4=\SI{2000}{\mega\becquerel}$. Three-quarters of what is made decays before it is used; that is
the price of an isotope short-lived enough to be safe.
\ans{\SI{2.0}{\giga\becquerel}}
\end{exbox}

\subsection{Positron Emission Tomography}

A tracer seen only through its gamma rays shows roughly where it is. Positron emission tomography (PET)
locates it far more precisely by using the annihilation of \chref{ch:decay}. Fluorine-18 ($\beta^+$,
$\Thalf=\SI{109.8}{min}$) is attached to a glucose-like molecule. Tissues that consume glucose quickly,
such as tumours and active brain regions, accumulate it. Each positron travels about a millimetre, meets
an electron and annihilates, $e^++e^-\to2\gamma$. The pair is almost at rest when it annihilates, so its
total momentum is nearly zero, and the two photons must leave in exactly opposite directions to keep it
zero (momentum conservation, \chref{ch:collisions}); their combined energy is the pair's rest energy,
$2\times\SI{0.511}{\MeV}$, so each carries \SI{511}{keV}.

A ring of detectors records pairs of photons that arrive at the same instant (Figure~\ref{fig:nucapp-pet}).
Each such pair defines a line through the point of annihilation, and a computer reconstructs the
distribution of the tracer from millions of these lines. The millimetre the positron travels before
annihilating sets a basic limit on how sharp the image can be.

\begin{figure}[htbp]
\centering
\tikzfig{nucapp_pet}
\figcaption{fig:nucapp-pet}{PET imaging. A positron from fluorine-18 in the tumour annihilates with an
electron. Notice that the two \SI{511}{keV} photons leave back to back, because the pair had almost no
momentum, and strike the ring at the same moment; the decay therefore lies somewhere on the dashed line
joining the two struck detectors.}
\end{figure}

\subsection{Therapy}

Ionising radiation kills cells by breaking their DNA, and rapidly dividing cancer cells are the least able
to repair it. Two strategies deliver the dose where it is wanted. The first places the source inside the
target. Iodine-131 ($\beta^-$, $\Thalf=\SI{8.02}{d}$), swallowed as sodium iodide, concentrates in the
thyroid gland, which uses iodine to make hormones, and its short-range $\beta^-$ particles destroy
overactive or cancerous thyroid tissue from within while sparing most of the body. The second aims beams
from outside: gamma rays from cobalt-60 or X-rays from an accelerator are directed at a tumour from many
directions, so that only the tumour, where the beams cross, receives the full dose.

An isotope inside a patient leaves by two routes at once. It decays, and the body also excretes it. If
excretion removes a fixed fraction of what is present per unit time (approximately true for many
substances), it is a second exponential process, characterised by a \emph{biological half-life}
$T_\text{bio}$, the time for the body to eliminate half of the substance by excretion alone. The physical
half-life is written $T_\text{phys}$.

\begin{derivbox}{Effective half-life}
Write $\lambda_\text{phys}=\ln2/T_\text{phys}$ and $\lambda_\text{bio}=\ln2/T_\text{bio}$ for the two
removal rates. The routes are independent, so the probabilities of removal per unit time add, exactly as
the partial decay constants of a branching decay add (\chref{ch:decaylaw}):
\[
\frac{dN}{dt}=-\lambda_\text{phys}N-\lambda_\text{bio}N=-(\lambda_\text{phys}+\lambda_\text{bio})N .
\]
The amount in the body therefore falls exponentially with $\lambda_\text{eff}=\lambda_\text{phys}
+\lambda_\text{bio}$. Dividing $\ln2$ by each side's half-life gives the result below.
\end{derivbox}

\begin{keybox}{Effective half-life}
For an isotope removed by physical decay and biological excretion together,
\[
\frac{1}{T_\text{eff}}=\frac{1}{T_\text{phys}}+\frac{1}{T_\text{bio}},
\]
and $T_\text{eff}$ is then used as an ordinary half-life. It is shorter than both, and close to the
shorter of the two when they are very different.
\end{keybox}

\begin{trapbox}{half-lives do not add}
It is the \emph{rates} that add, not the half-lives. Two removal routes make the substance disappear
faster, so $T_\text{eff}$ must be shorter than either. Adding half-lives, or averaging them, gives an
answer longer than one of them, which cannot be right.
\end{trapbox}

\section{Measuring Radiation Dose}\label{sec:nucapp-dose}

Using radiation safely requires a number for ``how much radiation a person has received''. Activity will
not do: it measures what a source emits, not what a body absorbs. The natural measure is the energy
deposited, per kilogram, since a small organ that absorbs a given energy is affected more than the whole
body absorbing the same energy. But energy alone is not enough either, and the reason shapes the
definitions.

\begin{defbox}{Absorbed dose and equivalent dose}
The \textbf{absorbed dose} $D$ is the energy $E$ deposited by ionising radiation per unit mass $m$ of
tissue,
\[
D=\frac{E}{m},
\]
measured in \textbf{grays}: $\SI{1}{\gray}=\SI{1}{J.kg^{-1}}$. The \textbf{equivalent dose} $H$ weights
the absorbed dose by a dimensionless \textbf{radiation weighting factor} $w_R$ that measures how damaging
each kind of radiation is for the same energy:
\[
H=w_RD,
\]
measured in \textbf{sieverts} (\si{\sievert}; dimensionally also \si{J.kg^{-1}}). The values used in
radiation protection are $w_R=1$ for gamma rays, X-rays and $\beta$ particles; $2$ for protons; $20$ for
$\alpha$ particles; and between about $2.5$ and $20$ for neutrons, depending on their energy.
\end{defbox}

Why should equal energies do unequal harm? An $\alpha$ particle deposits its few MeV in a track only
tens of micrometres long, so it ionises densely and breaks both strands of a DNA molecule in several nearby
places at once; such clustered damage is hard for the cell to repair. A gamma ray's electrons spread the
same energy thinly over centimetres, producing isolated breaks that are mostly repaired. The weighting
factor turns this difference into a number. Table~\ref{tab:nucapp-doses} puts the sievert on a human scale.

\begin{tablebox}{Typical equivalent doses}
\centering
\begin{tabular}{@{}lr@{}}\toprule
Exposure & Approximate equivalent dose\\\midrule
chest X-ray & \SI{0.02}{\milli\sievert}\\
transcontinental flight & \SI{0.04}{\milli\sievert}\\
world-average natural background, per year & \SI{2.4}{\milli\sievert}\\
CT scan of the abdomen & 8--\SI{10}{\milli\sievert}\\
occupational limit, per year (averaged over five years) & \SI{20}{\milli\sievert}\\
acute whole-body dose fatal to about half of those exposed & 4--\SI{5}{\sievert}\\\bottomrule
\end{tabular}\label{tab:nucapp-doses}
\end{tablebox}

Radiation harm is of two kinds. \emph{Deterministic} effects (radiation burns, cataracts, radiation
sickness) appear only above a threshold of roughly 0.5--\SI{1}{\sievert} received quickly, and grow more
severe with dose, because enough cells are killed to impair an organ. \emph{Stochastic} effects, chiefly
cancer arising from a cell that was damaged but survived, have no known threshold; what grows with dose is
the \emph{probability} of harm, not its severity. For radiation protection the risk is taken as about $5\%$
per sievert, so \SI{1}{\milli\sievert} adds roughly $0.05\times10^{-3}=1$ chance in $20\,000$.

That figure deserves candour. It comes from groups exposed to large doses (chiefly the survivors of
Hiroshima and Nagasaki) and is extended to small doses by assuming that risk is simply proportional to
dose, with no threshold. Below about \SI{100}{\milli\sievert} the extra cancers, if any, are too few to
detect against the natural cancer rate, so the straight-line assumption cannot be tested there. It is a
deliberately cautious rule for setting limits, not a measured law.

Everyone receives a natural dose. About half the world-average background comes from breathing radon-222,
a gas produced in the uranium-238 chain (\chref{ch:decay}) that seeps out of rock and soil, together
with its short-lived daughters. Parts of the Kerala coast, where the beach sand is rich in thorium-bearing
monazite, have natural background levels several times the world average.

\begin{exbox}{An internal $\alpha$ emitter}\label{ex:nucapp-po210dose}
A \SI{1.0}{\kilo\becquerel} source of polonium-210 ($\alpha$ emitter, $\Thalf=\SI{138.4}{d}$) lodges in
\SI{0.10}{kg} of tissue for one day. Each decay gives the $\alpha$ particle \SI{5.30}{\MeV}, all absorbed
in the tissue. Find the absorbed and equivalent doses.

\stepn{1} \textbf{Number of decays.} One day is short compared with the half-life (the activity falls by
only $0.5\%$), so the activity is effectively constant: decays $=\Act t=10^{3}\times86\,400=8.64\times10^{7}$.

\stepn{2} \textbf{Energy and absorbed dose.} $E=8.64\times10^{7}\times5.30\times1.602\times10^{-13}\
\si{J}=\SI{7.3e-5}{J}$, so $D=E/m=\SI{7.3e-4}{\gray}$.

\stepn{3} \textbf{Equivalent dose.} For $\alpha$ particles $w_R=20$: $H=20D=\SI{15}{\milli\sievert}$,
about six years of natural background delivered to that tissue in a single day. (The recoiling lead
nucleus also deposits its \SI{0.10}{\MeV} locally, which would add $2\%$.) Outside the body the same source
would be harmless, because its $\alpha$ particles cannot penetrate the dead outer layer of skin; inside, every
decay counts.
\ans{$D=\SI{0.73}{\milli\gray}$, $H\approx\SI{15}{\milli\sievert}$}
\end{exbox}

\section{Protection: Time, Distance and Shielding}\label{sec:nucapp-shielding}

Once dose can be measured, it can be reduced. For a source outside the body there are exactly three
handles: how long one stays near it, how far away one stands, and what one puts in between.

\begin{keybox}{Time, distance and shielding}
\textbf{Time:} at a steady dose rate, dose is proportional to exposure time.
\textbf{Distance:} from a small source, the dose rate falls as $1/r^2$ with distance $r$, so doubling the
distance quarters the dose.
\textbf{Shielding:} a gamma-ray beam's intensity $I$ falls exponentially with absorber thickness $x$,
\[
I=I_0e^{-\mu x}=I_0\,2^{-x/x_{1/2}},\qquad x_{1/2}=\frac{\ln2}{\mu},
\]
where $I_0$ is the intensity with no absorber, $\mu$ (in \si{cm^{-1}}) is the \textbf{linear
attenuation coefficient} and $x_{1/2}$ is the \textbf{half-value layer}, the thickness that halves the beam.
\end{keybox}

The distance rule is geometry: a small source sends its photons out equally in all directions, and at
distance $r$ they are spread over a sphere of area $4\pi r^2$, the same reasoning that gives the
inverse-square law of gravitation (\chref{ch:gravitation}). It assumes a source small compared with $r$ and
negligible absorption in the air between. The shielding rule needs a derivation.

\begin{derivbox}{Why gamma rays are attenuated exponentially}
A photon travels in a straight line until it interacts (by ejecting an electron, by scattering, or by
creating an electron--positron pair), and is then removed from the beam. In a thin layer of thickness
$dx$, each photon has the same small probability $\mu\,dx$ of interacting, whatever distance it has
already travelled. This is the memoryless rule of radioactive decay with distance in place of time. Of
$I$ photons entering the layer, $\mu I\,dx$ are removed:
\[
\frac{dI}{dx}=-\mu I\quad\Longrightarrow\quad I=I_0e^{-\mu x}.
\]
The thickness that halves the beam satisfies $e^{-\mu x_{1/2}}=\tfrac12$, so $x_{1/2}=\ln2/\mu$, the exact
analogue of $\Thalf=\ln2/\lambda$. The mean distance a photon travels before interacting, its mean free
path, is $1/\mu$, the analogue of the mean life $\tau=1/\lambda$ (and of the mean free path of a gas
molecule in \chref{ch:kinetic}). The coefficient $\mu$ grows with the density and atomic number of the
absorber and depends on the photon energy, which is why lead, dense and of high atomic number, is the
standard shield.
\end{derivbox}

Figure~\ref{fig:nucapp-shielding} sets the two laws side by side. Charged particles behave quite
differently. An $\alpha$ or $\beta$ particle loses its energy \emph{continuously}, in millions of small
ionising collisions, and stops after a definite range. That is why a sheet of paper stops every $\alpha$
particle from a source, whereas no thickness of lead stops every photon; it only reduces their number by a
known factor.

\begin{figure}[htbp]
\centering
\tikzfig{nucapp_shielding}
\figcaption{fig:nucapp-shielding}{Two ways to cut the dose from a gamma source. Left: each half-value layer
halves the transmitted intensity, so the curve has exactly the shape of radioactive decay and never
reaches zero. Right: doubling the distance from a point source quarters the dose rate; going from
\SI{1}{m} to \SI{4}{m} cuts it sixteenfold.}
\end{figure}

\begin{exbox}{How much lead?}\label{ex:nucapp-lead}
The half-value layer of lead for cobalt-60 gamma rays is about \SI{1.2}{cm}. What thickness reduces the
intensity by a factor of $1000$?

\stepn{1} \textbf{Number of half-value layers.} $2^{n}=1000$ gives $n=\log_21000=9.97$.

\stepn{2} \textbf{Thickness.} $x=n\,x_{1/2}=9.97\times1.2=\SI{12}{cm}$. (Check: ten half-value layers give
$2^{10}=1024$.)
\ans{about \SI{12}{cm}}
\end{exbox}

\begin{remarkbox}{Where the shielding formula stops working}
The exponential law describes a \emph{narrow} beam, where a photon that scatters is lost from the beam
for good. Behind a real shield exposed to a broad source, some scattered photons are deflected back
towards the person being protected, so more radiation gets through than $e^{-\mu x}$ predicts; engineers
multiply by a ``build-up factor'' greater than $1$. And since $\mu$ depends on photon energy, a half-value
layer quoted for cobalt-60 does not apply to technetium-99m, whose softer gamma rays are stopped far more
easily.
\end{remarkbox}

\section{Decay from Pure Randomness: a Monte Carlo Simulation}\label{sec:nucapp-montecarlo}

\chref{ch:decaylaw} began from a single assumption: each nucleus has the same fixed chance of decaying in
each small time interval, independently of its age and of every other nucleus. The exponential law was
derived from it with calculus. A computer can test the claim in a different way, by simply \emph{playing
out} the assumption with random numbers and seeing whether exponential decay emerges. A calculation that
uses random numbers to imitate a chance process is called a \emph{Monte Carlo simulation}, after the
casino.

Divide time into steps of length $\Delta t$. If $\lambda\Delta t$ is small, the probability that a given
nucleus decays during one step is $p=\lambda\Delta t$. Most computer languages provide a function that
returns a random number $u$ spread evenly between $0$ and $1$, so ``$u<p$'' happens with probability
exactly $p$: it is a die that comes up ``decay'' a fraction $p$ of the time.

\begin{methodbox}{Simulating radioactive decay}
\begin{enumerate}[label=\arabic*.,itemsep=1pt]
\item Choose $N_0$ (number of nuclei), $p=\lambda\Delta t$ and the number of steps. Set $N=N_0$.
\item In each step, for each of the $N$ surviving nuclei draw a random number $u$; count a decay if
      $u<p$. Subtract the number of decays from $N$ and record the new $N$.
\item Repeat for the required number of steps, and plot $N/N_0$ against the step number.
\item For lifetimes: take one nucleus, draw a random number each step until one falls below $p$, and record
      that step number $n$ as its lifetime. Repeat for many nuclei and make a histogram of $n$.
\end{enumerate}
A fixed starting value (\emph{seed}) for the random-number generator makes the run reproducible.
\end{methodbox}

Figure~\ref{fig:nucapp-montecarlo} shows the output for $p=0.05$ (that is, $\lambda=0.05$ per step with
$\Delta t=1$ step), three independent runs each for $N_0=1000$ and $N_0=100$, and $20\,000$ single-nucleus
lifetimes. No exponential was put in; the only rule was ``decay if $u<p$''. Yet the survival curves follow
$e^{-\lambda t}$, and the histogram of lifetimes follows the exponential distribution of lifetimes,
$\lambda e^{-\lambda t}$.

\begin{figure}[htbp]
\centering
\tikzfig{nucapp_montecarlo}
\figcaption{fig:nucapp-montecarlo}{A fixed-seed simulation with decay probability $p=0.05$ per step.
Left: survival curves for $N_0=1000$ (blue) and $N_0=100$ (orange), three runs each, against $e^{-\lambda t}$
(dashed). Notice that the large samples hug the curve while the small ones wander visibly. Right: histogram
of $20\,000$ simulated lifetimes, following $\lambda e^{-\lambda t}$; its mean, $20.05$ steps, reproduces the
mean life $\tau=1/\lambda=20$ steps.}
\end{figure}

The lifetime half of the simulation asks a different question: how long does one nucleus survive? It
survives each step with probability $1-p$, so the chance that it decays exactly at step $n$ is
\[
P(n)=(1-p)^{n-1}p,
\]
the \emph{geometric distribution}, which is the step-by-step version of $\lambda e^{-\lambda t}$.

\begin{derivbox}{The mean lifetime is $1/p$}
The mean of a distribution is the sum of each value times its probability (\chref{ch:toolkit},
Section~\ref{sec:averages}). Write $s=1-p$. Then
\[
\langle n\rangle=\sum_{n=1}^{\infty}n\,s^{n-1}p=p\,\frac{d}{ds}\sum_{n=0}^{\infty}s^{n}
=p\,\frac{d}{ds}\frac{1}{1-s}=\frac{p}{(1-s)^2}=\frac1p .
\]
With $p=\lambda\Delta t$ and $\Delta t=1$ step, $\langle n\rangle=1/\lambda$ steps: the mean life $\tau$,
exactly as found from calculus in \chref{ch:decaylaw}. For $p=0.05$ it is $20$ steps; the simulation's
$20\,000$ lifetimes average $20.05$.
\end{derivbox}

Three lessons emerge.
\begin{enumerate}[label=(\roman*)]
\item \textbf{The exponential law is a law of averages.} No individual nucleus follows it; the population
does. The number of decays in a step, among $N$ nuclei each with probability $p$, follows the
\emph{binomial distribution}, the distribution of the number of successes in $N$ independent trials. Its
mean is $Np$ and its standard deviation (typical spread about the mean) is $\sqrt{Np(1-p)}\approx\sqrt{Np}$.
The mean follows because each trial contributes $p$ on average; the spread is stated here without proof
(each trial contributes a variance $p(1-p)$, and variances of independent trials add).
\item \textbf{Fluctuations shrink as $1/\sqrt{N}$.} The relative scatter is $\sqrt{Np}/(Np)=1/\sqrt{Np}$.
For the $10^{20}$ or so nuclei of a real sample the curve is perfectly smooth; for $100$ nuclei it is
visibly jagged. Counting experiments inherit this: a measurement that records $n$ counts has an
uncertainty of about $\sqrt n$, so the fractional uncertainty $1/\sqrt n$ falls only slowly with counting
time (the uncertainties of \chref{ch:measurement} in their radioactive form).
\item \textbf{Small steps are required.} The rule $p=\lambda\Delta t$ is accurate only if $\lambda\Delta t\ll1$.
The exact per-step probability is $1-e^{-\lambda\Delta t}$, and for $\lambda\Delta t=0.05$ the two differ by
about $2.5\%$. This is the same small-step approximation that underlies every derivative.
\end{enumerate}

\begin{exbox}{Fluctuations in the first step}\label{ex:nucapp-mcspread}
In the simulation with $N_0=1000$ and $p=0.05$, what is the expected number of decays in the first step,
and its standard deviation?

\stepn{1} \textbf{Mean.} $Np=1000\times0.05=50$.

\stepn{2} \textbf{Spread.} $\sqrt{Np(1-p)}=\sqrt{1000\times0.05\times0.95}=6.9$. About two runs in three show
between $43$ and $57$ decays in the first step; the relative spread, $6.9/50\approx14\%$, would shrink to
$1.4\%$ with a hundred times as many nuclei.
\ans{$50\pm7$}
\end{exbox}

\section{Putting the Volume Together}\label{sec:nucapp-together}

The hardest examination problems, in JEE Advanced and the Physics Olympiad, rarely stay inside one
chapter. A single question may ask for an activity (\chref{ch:decaylaw}), turn it into a power with a
$Q$-value (\chref{ch:decay}), and then follow the product through a magnetic field or a detector. Such
problems are not harder in kind, only longer. The method is to break them into stages and give each stage
to the idea it belongs to.

\begin{methodbox}{Multi-concept problems}
\begin{enumerate}[label=\arabic*.,itemsep=1pt]
\item Split the problem into stages: \emph{energetics} (masses, $Q$-values, binding energies),
      \emph{kinetics} (the decay law, activity, chains, production), \emph{mechanics} (momentum, energy
      sharing, motion in fields) and \emph{counting} (detector efficiency, dose, statistics).
\item Solve each stage with the result it belongs to, carrying units through every step.
\item At every joint between stages, check units and the order of magnitude before going on.
\end{enumerate}
\end{methodbox}

\begin{exbox}{A nuclear battery}\label{ex:nucapp-battery}
Early space power sources, such as the heaters that kept the Soviet Lunokhod rovers warm through the lunar night, used the heat of polonium-210 ($\Thalf=\SI{138.4}{d}$,
molar mass \SI{210.0}{g.mol^{-1}}, $Q_\alpha=\SI{5.41}{\MeV}$). Find the thermal power of \SI{1.00}{g},
assuming the energy of every decay is absorbed, and the power one year later.

\stepn{1} \textbf{Kinetics: the activity.} $N=6.022\times10^{23}/210.0=2.87\times10^{21}$ nuclei and
$\lambda=0.693/(138.4\times86\,400\ \si{s})=\SI{5.80e-8}{s^{-1}}$, so $\Act=\lambda N=\SI{1.66e14}{\becquerel}$
(about \SI{4500}{\curie}).

\stepn{2} \textbf{Energetics: the power.} Both the $\alpha$ particle and the recoiling nucleus stop in the
source, so all of $Q$ becomes heat: $P=\Act Q=1.66\times10^{14}\times5.41\times1.602\times10^{-13}=\SI{144}{W}$.
One gram glows red hot.

\stepn{3} \textbf{Kinetics again: one year later.} One year is $365.25/138.4=2.64$ half-lives, so
$P=144\times2^{-2.64}=\SI{23}{W}$.

\stepn{4} \textbf{Check and interpret.} A loss of $84\%$ in a year is too fast for a long mission. Long missions
therefore use plutonium-238 ($\Thalf=\SI{88}{\year}$): far less power per gram, but almost all of it is
still there after a decade. Choosing an isotope is a trade between power and lifetime, because both are set
by the same $\lambda$.
\ans{\SI{144}{W}; \SI{23}{W}}
\end{exbox}

Examinations also test understanding with \emph{assertion--reason} items: two statements, an Assertion (A)
and a Reason (R), with four standard options. (A) Both are true and R correctly explains A. (B) Both are
true but R does not explain A. (C) A is true but R is false. (D) A is false but R is true. The method is to
judge each statement on its own first, and only then ask whether R is the actual cause of A, not merely a
true statement on the same topic. Eight such items, drawn from all five chapters, are among the Problem
Types below.

\section{Summary}\label{sec:summary-nucapps}

The exponential law of \chref{ch:decaylaw} does practical work in four ways. Read backwards, it dates a
sample from one measured ratio. Applied inside the body, together with excretion, it schedules medical
isotopes and limits the dose they give. Applied to photons crossing matter, the same memoryless rule gives
exponential shielding. And it emerges unaided from random numbers, as a law of averages whose fluctuations
shrink as $1/\sqrt N$.

\begin{itemize}
\item Dating a closed sample with no initial daughter: $t=\Thalf\log_2(1+N_Y/N_X)$; with a known start,
      $t=\Thalf\log_2(\Act_0/\Act)$. A branching parent with daughter fraction $f$:
      $N_Y/N_X=f(e^{\lambda t}-1)$, $\lambda$ the total decay constant.
\item Radiocarbon: living carbon has about $13.8$ decays per minute per gram ($\Thalf=\SI{5730}{\year}$);
      useful to about \SI{50000}{\year}. Rocks use uranium--lead, potassium--argon and rubidium--strontium.
\item Medicine: tracers (technetium-99m), PET (two back-to-back \SI{511}{keV} photons), therapy
      (iodine-131, cobalt-60). Dispatch activity $\Act_\text{use}2^{t/\Thalf}$;
      $1/T_\text{eff}=1/T_\text{phys}+1/T_\text{bio}$.
\item Dose: $D=E/m$ in grays; $H=w_RD$ in sieverts ($w_R=1$ for $\gamma$, $\beta$; $20$ for $\alpha$).
      Deterministic effects have a threshold; stochastic risk is taken as about $5\%$ per sievert.
\item Protection: dose $\propto$ time, $\propto1/r^2$, and $I=I_0e^{-\mu x}=I_02^{-x/x_{1/2}}$ with
      $x_{1/2}=\ln2/\mu$ (narrow beam). Charged particles have a definite range; photons do not.
\item Simulation: decay with probability $p=\lambda\Delta t$ per step reproduces $e^{-\lambda t}$; lifetimes are
      geometric with mean $1/p$; decays per step are binomial with spread $\sqrt{Np(1-p)}$; a count $n$ has
      uncertainty $\sqrt n$.
\end{itemize}

Because this chapter closes the volume, the results of all five chapters are gathered in the
formula sheet of Table~\ref{tab:nucapp-formulas}, followed by the ten errors that cost the most marks.

\begin{tablebox}{Formula sheet for the volume}
\renewcommand{\arraystretch}{1.22}
\begin{tabular}{@{}p{0.29\linewidth}p{0.67\linewidth}@{}}\toprule
Quantity & Relation\\\midrule
Nuclear radius, density & $R=R_0A^{1/3}$, $R_0\approx\SI{1.2}{\fm}$; $\rho\approx\SI{2.3e17}{kg.m^{-3}}$ for all nuclei\\
Mass--energy & $1\,\si{\amu}\times c^2=\SI{931.5}{\MeV}$\\
Binding energy & $\Eb=[ZM(\iso{1}{H})+(A-Z)m_n-M]c^2$ (atomic masses)\\
Energy from rearrangement & $Q=\sum E_{\mathrm{b},\text{products}}-\sum E_{\mathrm{b},\text{reactants}}$\\
Reaction X$(a,b)$Y & $Q=[(m_a+m_X)-(m_b+m_Y)]c^2$; threshold $K_\text{th}=|Q|(1+m_a/m_X)$\\
$\alpha$ decay & $Q=[M_X-M_Y-M_\text{He}]c^2$; $K_\alpha=Q(A-4)/A$\\
$\beta^-$ decay, electron capture & $Q=[M_X-M_Y]c^2$\\
$\beta^+$ decay & $Q=[M_X-M_Y-2m_e]c^2$ (needs $\Delta M\,c^2>\SI{1.022}{\MeV}$)\\
Counting decays in a series & $n_\alpha=\Delta A/4$; $n_\beta=Z_f-Z_i+2n_\alpha$\\
Decay law & $N=N_0e^{-\lambda t}=N_02^{-t/\Thalf}$; $\Act=\lambda N$\\
Half-life, mean life & $\Thalf=\ln2/\lambda=0.693/\lambda$; $\tau=1/\lambda=1.443\,\Thalf$\\
Branching, effective half-life & $\lambda=\sum_i\lambda_i$; $1/\Thalf=\sum_i1/T_i$; $1/T_\text{eff}=1/T_\text{phys}+1/T_\text{bio}$\\
Daughter in a chain & $N_Y=\dfrac{\lambda_XN_{X,0}}{\lambda_Y-\lambda_X}\left(e^{-\lambda_Xt}-e^{-\lambda_Yt}\right)$;
  $t_{\max}=\dfrac{\ln(\lambda_Y/\lambda_X)}{\lambda_Y-\lambda_X}$\\
Secular equilibrium & $\lambda_XN_X=\lambda_YN_Y$\\
Constant production & $\Act=\Prod(1-e^{-\lambda t})$, saturating at $\Prod$\\
Dating & $t=\Thalf\log_2(1+N_Y/N_X)$, or $t=\Thalf\log_2(\Act_0/\Act)$\\
Moderation (head-on) & $\phi=K'/K=\big((A-1)/(A+1)\big)^2$\\
Chain reaction & $k^n$ neutrons per starting neutron after $n$ generations\\
Coulomb barrier & $U_B=\dfrac{1}{4\pi\varepsilon_0}\dfrac{q_1q_2}{r}$, with $\dfrac{e^2}{4\pi\varepsilon_0}=\SI{1.44}{\MeV.\fm}$\\
Dose & $D=E/m$ (\si{\gray}); $H=w_RD$ (\si{\sievert})\\
Shielding, distance & $I=I_0e^{-\mu x}=I_02^{-x/x_{1/2}}$; dose rate $\propto1/r^2$\\\bottomrule
\end{tabular}\label{tab:nucapp-formulas}
\end{tablebox}

\subsection*{Ten errors to avoid}
\begin{enumerate}[label=\arabic*.,itemsep=1pt]
\item Using the bare proton mass with atomic masses. Use $M(\iso{1}{H})$ so that the electron masses cancel.
\item Forgetting the $2m_ec^2$ in $\beta^+$ $Q$-values, or wrongly adding it in $\beta^-$ decay and electron
      capture.
\item Giving all of $Q$ to the $\alpha$ particle. The recoiling daughter takes the fraction $4/A$.
\item Treating the $\beta$ endpoint energy as the typical electron energy. The average is about $Q/3$.
\item Leaving $\lambda$ in \si{\year^{-1}} when computing an activity in becquerels.
\item Thinking a sample is ``gone'' after two half-lives. A quarter remains.
\item Confusing ``fraction remaining'' with ``fraction decayed''.
\item Using partial half-lives as if each channel ran alone. The sample decays with the total $\lambda$.
\item Forgetting to subtract the background count before finding a half-life from counts.
\item Confusing grays (energy per kilogram) with sieverts (biological effect).
\end{enumerate}

\section{Problem Types}\label{sec:types-nucapps}

Sixteen types follow. The first eight belong to this chapter: dating, medical activity, dose, shielding,
counting statistics, and a multi-stage problem. Nearly all of them are the decay law in disguise, so the
first question is always which quantity is decaying and what plays the part of the half-life. The last
eight are assertion--reason items (options as defined in Section~\ref{sec:nucapp-together}) that revisit the
whole volume; for these, judge each statement alone before asking whether the reason is the cause.

\archq{Radiocarbon dating}
{``The carbon-14 activity per gram of an old sample is a given fraction of the living value; find its age.''}
{Put both activities on the same basis (per gram of carbon), take the ratio
$\Act_\text{living}/\Act_\text{sample}$, and use $t=\Thalf\log_2(\text{ratio})$ with $\Thalf=\SI{5730}{\year}$.
A ratio that is a power of $2$ is a whole number of half-lives.}
{Treating decay as linear: losing seven-eighths does not take seven-eighths of some fixed time. Or
comparing a total activity with a per-gram activity.}
{The activity per gram of carbon in an old bone is one-eighth of that in living bone ($\Thalf=\SI{5730}{\year}$).
The age of the bone is about:}
{\SI{5730}{\year}}{\SI{11460}{\year}}{\SI{17190}{\year}}{\SI{45840}{\year}\archans{C}{A factor of $8=2^3$ is
three half-lives, $3\times5730=\SI{17190}{\year}$. Option (A) is one half-life and (B) two (a factor of $2$ or
$4$); (D) multiplies the half-life by $8$, treating the factor as a number of half-lives.}}

\archq{Parent--daughter dating of rocks}
{``A mineral contains so many daughter atoms per parent atom; find its age.''}
{Assume a closed system with no initial daughter and use $t=\Thalf\log_2(1+N_Y/N_X)$. The $1+$ restores the
original number of parents, $N_X+N_Y$.}
{Writing $\log_2(N_Y/N_X)$, which forgets that the original parent count is $N_X+N_Y$; or mixing $\ln$ with
$\log_2$.}
{A rock contains three atoms of a stable daughter for every atom of its parent ($\Thalf=\SI{1.3e9}{\year}$),
and had no daughter when it formed. Its age is:}
{\SI{3.9e9}{\year}}{\SI{2.6e9}{\year}}{\SI{2.1e9}{\year}}{\SI{1.8e9}{\year}\archans{B}{$1+N_Y/N_X=4=2^2$,
so $t=2\Thalf=\SI{2.6e9}{\year}$. Option (A) multiplies the half-life by the ratio $3$; (C) is
$\Thalf\log_23$, forgetting the $1+$; (D) is $\Thalf\ln4$, using the natural logarithm where the base-$2$
logarithm belongs.}}

\archq{Dating with a branching parent}
{``Only part of the parent's decays produce the measured daughter'' (potassium--argon).}
{The parent disappears with its \emph{total} $\lambda$, so decayed parents number $N_X(e^{\lambda t}-1)$, and only
the fraction $f$ became the measured daughter: $N_Y/N_X=f(e^{\lambda t}-1)$, so $t=\Thalf\log_2(1+N_Y/fN_X)$.}
{Ignoring the branching (using $1+N_Y/N_X$), or slowing the clock by $1/f$ as though the parent decayed only
through the measured channel.}
{Potassium-40 ($\Thalf=\SI{1.25e9}{\year}$) decays to argon-40 in $10.7\%$ of its decays. A lava sample that
lost all its argon on solidifying now has $0.107$ argon-40 atoms per potassium-40 atom. Its age is about:}
{\SI{0.18e9}{\year}}{\SI{0.87e9}{\year}}{\SI{1.25e9}{\year}}{\SI{11.7e9}{\year}\archans{C}{$1+N_Y/(fN_X)=1+0.107/0.107
=2$, so exactly one half-life, \SI{1.25e9}{\year}. Option (A) ignores the branching, $\Thalf\log_2(1.107)$;
(B) is $\Thalf\ln2$, mixing logarithms; (D) divides the age by $f$, as if the parent decayed only by the
argon channel, which also exceeds the age of the universe.}}

\archq{Medical activity and effective half-life}
{``Activity needed at dispatch for a given activity at use,'' or ``an isotope eliminated both physically and
biologically.''}
{Work backwards with $\Act_\text{dispatch}=\Act_\text{use}\times2^{t/\Thalf}$. For combined removal the rates
add: $1/T_\text{eff}=1/T_\text{phys}+1/T_\text{bio}$; then use $T_\text{eff}$ as an ordinary half-life.}
{Adding or averaging the half-lives. Two removal routes act faster than either alone, so $T_\text{eff}$ is
shorter than both.}
{An isotope has a physical half-life of \SI{6}{h} and a biological half-life of \SI{12}{h}. Its effective
half-life in the body is:}
{\SI{4}{h}}{\SI{6}{h}}{\SI{9}{h}}{\SI{18}{h}\archans{A}{$1/T_\text{eff}=1/6+1/12=3/12$, so
$T_\text{eff}=\SI{4}{h}$, shorter than both. Option (B) keeps only the shorter half-life; (C) averages and (D)
adds the half-lives, both longer than the physical half-life alone, which is impossible.}}

\archq{Absorbed and equivalent dose}
{``Energy deposited in a mass of tissue,'' grays versus sieverts, or an internal emitter.}
{$D=E/m$ in grays; $H=w_RD$ in sieverts, with $w_R=20$ for $\alpha$ and $1$ for $\gamma$, $\beta$. For an
internal emitter, $E=(\text{number of decays})\times(\text{energy per decay})$, with decays $\approx\Act t$
when $t\ll\Thalf$.}
{Forgetting $w_R$ (quoting grays as sieverts), dividing by it instead of multiplying, or multiplying by the
mass instead of dividing.}
{A \SI{2.0}{kg} organ absorbs \SI{0.020}{J} of energy from $\alpha$ particles ($w_R=20$). The equivalent dose is:}
{\SI{0.5}{\milli\sievert}}{\SI{200}{\milli\sievert}}{\SI{10}{\milli\sievert}}{\SI{800}{\milli\sievert}
\archans{B}{$D=0.020/2.0=\SI{0.010}{\gray}$ and $H=20\times0.010=\SI{0.20}{\sievert}=\SI{200}{\milli\sievert}$.
Option (A) divides by $w_R$; (C) forgets it, quoting the gray value as sieverts; (D) multiplies by the mass
instead of dividing.}}

\archq{Distance and shielding together}
{``Dose rate at a new distance, behind an absorber.''}
{Two independent factors multiply: $(r_1/r_2)^2$ for distance and $2^{-x/x_{1/2}}$ for shielding.}
{Using $1/r$ instead of $1/r^2$, or applying only one of the two factors.}
{The dose rate \SI{1.0}{m} from a gamma source is \SI{64}{\micro\sievert/h}. At \SI{2.0}{m}, behind lead two
half-value layers thick, the dose rate is:}
{\SI{16}{\micro\sievert/h}}{\SI{8}{\micro\sievert/h}}{\SI{4}{\micro\sievert/h}}{\SI{32}{\micro\sievert/h}
\archans{C}{Distance: $(1/2)^2=1/4$; shielding: $2^{-2}=1/4$; so $64/16=\SI{4}{\micro\sievert/h}$. Option (A)
applies only one factor; (B) uses $1/r$ for distance; (D) treats two half-value layers as a single halving and
ignores the distance.}}

\archq{Counting statistics}
{``How uncertain is a count?'' or ``how long must one count for a given precision?''}
{A count of $n$ random events has uncertainty $\sqrt n$, so a fractional uncertainty of $1/\sqrt n$. At a
steady rate, $n$ grows in proportion to counting time.}
{Taking the fractional uncertainty as $1/n$, so that precision seems to improve in proportion to time.}
{To halve the fractional uncertainty of a measured count rate, the counting time must be multiplied by:}
{$\sqrt2$}{$2$}{$8$}{$4$\archans{D}{Fractional uncertainty $=1/\sqrt n$; halving it needs $4n$ counts, so four
times the time. Option (B) assumes uncertainty $\propto1/n$; (A) inverts the square root; (C) cubes instead of
squaring.}}

\archq{Multi-concept: power from a decaying source}
{Several stages in one problem: activity, then energy per decay, then time dependence.}
{Activity $\Act=\lambda N$; power $P=\Act Q$ when all the decay energy is absorbed; and $P$ decays with the
same half-life as $\Act$, so after $t$, $P=P_0\,2^{-t/\Thalf}$.}
{Letting the power fall linearly with time, or treating a source as exhausted after a couple of half-lives.}
{A curium-242 heat source ($\Thalf=\SI{163}{d}$), of the kind tested for early space missions, delivers
\SI{120}{W} at launch. One year ($365$ days) later it delivers about:}
{\SI{0}{W}}{\SI{25}{W}}{\SI{54}{W}}{\SI{60}{W}\archans{B}{One year is $365/163=2.24$ half-lives, so
$P=120\times2^{-2.24}=\SI{25}{W}$. Option (A) treats the source as used up; (C) divides the power by the
number of half-lives, $120/2.24$; (D) halves it only once.}}

\archq{Assertion--reason: the half-life and the amount present}
{An assertion about half-life with a reason about the order of the decay process.}
{A first-order law, $dN/dt=-\lambda N$, gives $N=N_0e^{-\lambda t}$, and the time to halve is $\ln2/\lambda$
whatever $N_0$ is.}
{Accepting a true R as an explanation without checking that it actually implies A; here it does.}
{A: The half-life of a radioactive sample does not depend on how much of it is present.\quad R: Radioactive
decay is a first-order process.}
{both true; R explains A}{both true; R does not explain A}{A true; R false}{A false; R true\archans{A}{Both
statements are true, and R is the reason: for a first-order process $\Thalf=0.693/\lambda$, independent of
$N_0$. A zeroth- or second-order process would have a half-life that depends on the amount.}}

\archq{Assertion--reason: the energy of $\alpha$ particles}
{An assertion about the spread of particle energies with a reason about the number of final bodies.}
{In a two-body decay, momentum and energy conservation fix both kinetic energies uniquely
(\chref{ch:decay}). Three bodies can share the energy in many ways.}
{Confusing the sharp $\alpha$ lines with the continuous $\beta$ spectrum, which comes from three-body decay.}
{A: All $\alpha$ particles emitted in a given transition have the same kinetic energy.\quad R: $\alpha$ decay
is a two-body decay.}
{both true; R explains A}{both true; R does not explain A}{A true; R false}{A false; R true\archans{A}{Two
bodies with equal and opposite momenta and a fixed total energy $Q$ can share it in only one way,
$K_\alpha=Q(A-4)/A$, so R explains A.}}

\archq{Assertion--reason: neutron excess in heavy nuclei}
{An assertion about the neutron excess with a reason that is true but beside the point.}
{Ask \emph{why} heavy nuclei need extra neutrons: protons repel each other by the long-range Coulomb force,
which grows as $Z^2$, while extra neutrons add strong-force attraction without repulsion (\chref{ch:nucleus}).}
{Accepting a true fact about the neutron as the explanation because it mentions neutrons.}
{A: Heavy stable nuclei contain more neutrons than protons.\quad R: The neutron is slightly heavier than the
proton.}
{both true; R explains A}{both true; R does not explain A}{A true; R false}{A false; R true\archans{B}{Both are
true, but the neutron excess comes from Coulomb repulsion among the protons, not from the neutron's mass
(which would, if anything, favour protons, since a free neutron decays into a proton).}}

\archq{Assertion--reason: activity halves each half-life}
{An assertion about activity with a reason that misstates what the half-life depends on.}
{$\Act=\lambda N$, so activity falls in step with $N$ and halves every $\Thalf$; $\Thalf=\ln2/\lambda$ depends
only on the nuclide.}
{Agreeing with R because larger samples have larger activities: the activity depends on $N$, the half-life
does not.}
{A: The activity of a pure sample of one radioactive nuclide halves every half-life.\quad R: The half-life of a
nuclide depends on the number of nuclei initially present.}
{both true; R explains A}{both true; R does not explain A}{A true; R false}{A false; R true\archans{C}{A is
true because $\Act=\lambda N$ and $N$ halves every half-life. R is false: $\Thalf=\ln2/\lambda$ is fixed by the
nuclide alone.}}

\archq{Assertion--reason: fusing iron}
{An assertion about energy release with a reason about the binding-energy curve.}
{Energy is released when the products have a higher binding energy per nucleon $\Eb/A$ than the reactants.
Near the peak at iron, any rearrangement lowers $\Eb/A$ (\chref{ch:reactions}).}
{Thinking every fusion releases energy. Only fusion of nuclei lighter than iron does.}
{A: Fusion of two $\iso{56}{Fe}$ nuclei releases energy.\quad R: $\iso{56}{Fe}$ lies near the maximum of the
$\Eb/A$ curve.}
{both true; R explains A}{both true; R does not explain A}{A true; R false}{A false; R true\archans{D}{Fusing two
iron nuclei would make a nucleus with \emph{lower} $\Eb/A$, absorbing energy, so A is false; R is true, and is
precisely why A fails.}}

\archq{Assertion--reason: where the $\beta$ electron comes from}
{An assertion that the emitted electron pre-existed in the nucleus.}
{In $\beta^-$ decay a neutron changes into a proton, $n\to p+e^-+\bar\nu_e$; the electron is \emph{created} in
the decay (\chref{ch:decay}). The uncertainty principle forbids an electron confined to a nucleus.}
{Reasoning that because something comes out, it must have been inside.}
{A: The electron emitted in $\beta^-$ decay was present inside the nucleus before the decay.\quad R: In $\beta^-$
decay the atomic number increases by one.}
{both true; R explains A}{both true; R does not explain A}{A true; R false}{A false; R true\archans{D}{A is
false: the electron is created when a neutron becomes a proton. R is true, since one more proton means $Z$
rises by one.}}

\archq{Assertion--reason: mean life and half-life}
{An assertion comparing $\tau$ and $\Thalf$ with a reason giving the relation between them.}
{$\Thalf=\tau\ln2$, so $\tau=\Thalf/\ln2=1.443\,\Thalf$.}
{Thinking the average lifetime must be the time by which half have decayed. The few long-lived survivors pull
the mean above the median.}
{A: The mean life of a radioactive nucleus is longer than its half-life.\quad R: $\tau=\Thalf/\ln2$ and $\ln2<1$.}
{both true; R explains A}{both true; R does not explain A}{A true; R false}{A false; R true\archans{A}{Both are
true, and dividing by $\ln2=0.693<1$ makes $\tau$ larger than $\Thalf$, so R is exactly the reason.}}

\archq{Assertion--reason: can lead stop every gamma ray?}
{An assertion about complete stopping with a reason about exponential attenuation.}
{An exponential $I_0e^{-\mu x}$ never reaches zero: every half-value layer halves what is left. Charged particles,
which have a definite range, are different (Section~\ref{sec:nucapp-shielding}).}
{Carrying over the idea of a range from $\alpha$ and $\beta$ particles to photons.}
{A: Gamma radiation from a source can be completely stopped by a sufficiently thick lead sheet.\quad R:
Gamma-ray intensity decreases exponentially with thickness.}
{both true; R explains A}{both true; R does not explain A}{A true; R false}{A false; R true\archans{D}{R is true,
and it shows A is false: exponential attenuation reduces the intensity by any desired factor but never to zero.}}

\section{Practice Problems}\label{sec:practice-nucapps}

The problems below are of standard examination difficulty and follow the order of the chapter. Answers, with
brief working, follow the list. Use $\SI{1}{\year}=\SI{3.156e7}{s}$ and
$\SI{1}{\MeV}=\SI{1.602e-13}{J}$.

\begin{enumerate}[itemsep=3pt]
\item\label{pp:nucapp-bone} A bone retains $60\%$ of the carbon-14 activity per gram of carbon of living bone.
      Find its age.
\item\label{pp:nucapp-wood} A \SI{5.0}{g} sample of carbon from a wooden tool shows $25$ decays per minute.
      Living carbon shows $13.8$ decays per minute per gram. Find the tool's age.
\item\label{pp:nucapp-dino} Explain why radiocarbon cannot date a dinosaur bone (66 million years old).
      What fraction of the carbon-14 would remain?
\item\label{pp:nucapp-kar} Potassium-40 ($\Thalf=\SI{1.248e9}{\year}$) decays to argon-40 in $10.72\%$ of its
      decays (by electron capture) and otherwise to calcium-40. A volcanic rock has $0.100$ argon-40 atoms per
      potassium-40 atom. Find its age, assuming all argon escaped while the rock was molten.
\item\label{pp:nucapp-rbsr} A meteorite contains $0.065$ atoms of radiogenic strontium-87 per atom of
      rubidium-87 ($\Thalf=\SI{4.97e10}{\year}$). Find its age.
\item\label{pp:nucapp-equal} A rock contains equal numbers of parent and daughter atoms. How old is it, in
      half-lives? What if $N_Y/N_X=15$?
\item\label{pp:nucapp-i131} Iodine-131 has $T_\text{phys}=\SI{8.02}{d}$. Suppose the thyroid eliminates iodine
      with $T_\text{bio}=\SI{80}{d}$. Find the effective half-life.
\item\label{pp:nucapp-f18} Fluorine-18 ($\Thalf=\SI{109.8}{min}$) is made at a cyclotron \SI{2.0}{h} by road
      from a PET centre that needs \SI{400}{\mega\becquerel} on arrival. What activity must leave the cyclotron?
\item\label{pp:nucapp-effshort} Show that $T_\text{eff}$ is always shorter than both $T_\text{phys}$ and
      $T_\text{bio}$, and that it is close to the shorter of the two when they are very different.
\item\label{pp:nucapp-gamma70} A \SI{70}{kg} worker absorbs \SI{0.35}{J} of gamma-ray energy spread through
      the body. Find $D$ and $H$. What would $H$ be if the same energy were delivered by $\alpha$ particles?
\item\label{pp:nucapp-lethal} A whole-body dose of \SI{5}{\gray} is likely to be fatal. How much energy is this
      for a \SI{70}{kg} person, and by how much would it raise the body's temperature (specific heat of tissue
      $s\approx\SI{4186}{J.kg^{-1}.K^{-1}}$, as in \chref{ch:thermal})? What does this say about how radiation
      harms the body?
\item\label{pp:nucapp-invsq} A technician moves from \SI{0.50}{m} to \SI{2.0}{m} from a point source. By what
      factor does the dose rate change?
\item\label{pp:nucapp-worker} The dose rate \SI{1.0}{m} from a source is \SI{80}{\micro\sievert/h}. Find the
      rate at \SI{4.0}{m} behind \SI{2.4}{cm} of lead (half-value layer \SI{1.2}{cm}).
\item\label{pp:nucapp-mu} Find the linear attenuation coefficient $\mu$ of lead for cobalt-60 gamma rays
      (half-value layer \SI{1.2}{cm}), and the fraction of a narrow beam transmitted by \SI{5.0}{cm}.
\item\label{pp:nucapp-pexact} For $\lambda\Delta t=0.05$, compute the exact per-step decay probability and the
      percentage error of using $p=\lambda\Delta t$.
\item\label{pp:nucapp-count} A detector records $400$ counts in a minute. Estimate the uncertainty in the count
      rate. How long must one count to measure the rate to $1\%$?
\item\label{pp:nucapp-chainsim} Extend the simulation method of Section~\ref{sec:nucapp-montecarlo} to a chain
      $\mathrm X\to\mathrm Y\to{}$stable, with decay probabilities $p_X=0.02$ and $p_Y=0.04$ per step and no Y at
      the start. Describe the algorithm, state the shape the curve for Y should have, and find the step at
      which it should peak, using the time of maximum daughter population from \chref{ch:decaylaw},
      $t_{\max}=\ln(\lambda_Y/\lambda_X)/(\lambda_Y-\lambda_X)$.
\end{enumerate}

\subsection*{Answers}
{\small
\begin{enumerate}[itemsep=2pt]
\item $t=5730\log_2(1/0.60)=5730\times0.737=\SI{4.2e3}{\year}$.
\item The living value for \SI{5.0}{g} is $5.0\times13.8=69$ per minute; ratio $69/25=2.76$;
      $t=5730\log_22.76=5730\times1.465=\SI{8.4e3}{\year}$.
\item $66\times10^{6}/5730\approx11\,500$ half-lives; fraction $2^{-11\,500}\approx10^{-3467}$. Not a single
      atom of the original carbon-14 remains (a living ten-tonne animal holds only about $10^{17}$ of them), so
      the clock ran out after about \SI{50000}{\year}. Dinosaurs are dated by the rocks around them, with
      long-lived parents.
\item With $f=0.1072$: $0.100=0.1072\,(e^{\lambda t}-1)$, so $e^{\lambda t}=1.933$ and
      $t=1.248\times10^{9}\times\ln1.933/0.693=\SI{1.19e9}{\year}$.
\item $t=4.97\times10^{10}\log_2(1.065)=4.97\times10^{10}\times0.0909=\SI{4.5e9}{\year}$, the age of the solar
      system.
\item $N_Y/N_X=1$: $\log_22=1$ half-life. $N_Y/N_X=15$: $\log_216=4$ half-lives.
\item $1/T_\text{eff}=1/8.02+1/80$, so $T_\text{eff}=\SI{7.3}{d}$: almost the physical half-life, because
      excretion is ten times slower.
\item $\SI{2.0}{h}=\SI{120}{min}$ is $120/109.8=1.093$ half-lives, so $400\times2^{1.093}=\SI{853}{\mega\becquerel}$.
\item $1/T_\text{eff}=1/T_\text{phys}+1/T_\text{bio}$ exceeds both $1/T_\text{phys}$ and $1/T_\text{bio}$, so
      $T_\text{eff}$ is smaller than both. If $T_\text{bio}\gg T_\text{phys}$, then $1/T_\text{bio}$ is negligible
      and $T_\text{eff}\approx T_\text{phys}$ (and vice versa).
\item $D=0.35/70=\SI{5.0}{\milli\gray}$; $H=\SI{5.0}{\milli\sievert}$ ($w_R=1$). As $\alpha$ particles,
      $H=20\times5.0=\SI{100}{\milli\sievert}$.
\item $E=5\times70=\SI{350}{J}$; $\Delta T=E/(ms)=350/(70\times4186)=\SI{0.0012}{K}$. A lethal dose would not
      noticeably warm the body. The harm comes not from heat but from ionisation: the energy arrives in
      concentrated packets that break chemical bonds in DNA at particular points.
\item $(0.50/2.0)^2=1/16$: the dose rate falls sixteenfold.
\item Distance: $(1/4)^2=1/16$; shielding: two half-value layers, $1/4$. Rate $=80/64=\SI{1.25}{\micro\sievert/h}$.
\item $\mu=0.693/1.2=\SI{0.58}{cm^{-1}}$; transmitted fraction $e^{-0.58\times5.0}=0.056$, that is, $5.6\%$.
\item $1-e^{-0.05}=0.0488$; using $0.05$ overestimates it by $2.5\%$.
\item $\sqrt{400}=20$, so the rate is $400\pm20$ per minute, a $5\%$ uncertainty. For $1\%$ one needs
      $1/\sqrt n=0.01$, that is, $n=10^4$ counts, which takes $10^4/400=\SI{25}{min}$.
\item In each step, decide for each X nucleus whether it decays (probability $p_X$) and for each Y nucleus
      whether it decays (probability $p_Y$); subtract the X decays from $N_X$, add them to $N_Y$, and subtract
      the Y decays from $N_Y$. $N_Y$ rises from zero, peaks and falls, the shape of a daughter whose parent
      outlives it only moderately. With $\lambda\approx p$ per step, $t_{\max}=\ln(0.04/0.02)/(0.04-0.02)=34.7$,
      so the peak is near step $35$ (step $34$ if the exact per-step probabilities $1-e^{-\lambda}$ are used).
\end{enumerate}
}

\section{Challenge Problems}\label{sec:challenge-nucapps}

The five problems below each combine ideas from several chapters of the volume, in the style of JEE
Advanced and Olympiad papers. Several are drawn from historical experiments. Break each into stages as
in Section~\ref{sec:nucapp-together}.

\begin{challengebox}{The $\alpha$ line of americium-241 (\chsref{ch:nucleus}{ch:decay})}
The smoke-detector isotope americium-241 (\chref{ch:decaylaw}) decays by $\alpha$ emission to neptunium-237.
The atomic masses are $M(\iso{241}{Am})=\SI{241.056829}{\amu}$, $M(\iso{237}{Np})=\SI{237.048173}{\amu}$ and
$M(\iso{4}{He})=\SI{4.002603}{\amu}$. (a) Find $Q_\alpha$ and the $\alpha$ energy expected for a decay to the
ground state of neptunium-237. (b) The strongest $\alpha$ line is measured at \SI{5.486}{\MeV}. Show that this
decay leaves the neptunium in an excited state, and find its excitation energy; what is then emitted? (c) Find
the radius of the path of a \SI{5.486}{\MeV} $\alpha$ particle (mass \SI{6.645e-27}{kg}, charge $2e$) moving
perpendicular to a magnetic field of \SI{0.50}{T}.
\challengeans{(a) $Q_\alpha=[241.056829-237.048173-4.002603](931.494)=\SI{5.638}{\MeV}$, and a decay to the
ground state would give $K_\alpha=Q\times237/241=\SI{5.545}{\MeV}$. (b) Working backwards, the observed line
corresponds to $Q'=K_\alpha A/(A-4)=5.486\times241/237=\SI{5.579}{\MeV}$, less than $Q_\alpha$ by
\SI{0.060}{\MeV}. The missing \SI{60}{\keV} is left as excitation of the daughter, which drops to its ground
state by emitting a gamma ray: the measured \SI{59.5}{\keV} line of americium sources. The two-body sharing
rule of \chref{ch:decay} thus turns an $\alpha$ energy into a map of the daughter's levels. (c)
$p=\sqrt{2mK}=\sqrt{2\times6.645\times10^{-27}\times5.486\times1.602\times10^{-13}}=\SI{1.08e-19}{kg.m.s^{-1}}$, so
$r=p/(2eB)=1.08\times10^{-19}/(2\times1.602\times10^{-19}\times0.50)=\SI{0.67}{m}$. A field of this size
bends $\alpha$ particles only gently, which is why magnetic analysis of $\alpha$ lines needs strong fields
and long flight paths.}
\end{challengebox}

\begin{challengebox}{Activation and counting (\chref{ch:decaylaw})}
A target is irradiated so that an isotope X ($\Thalf=\SI{2.0}{h}$) is produced at a constant
$1.0\times10^{8}$ nuclei per second for \SI{4.0}{h}. Counting starts \SI{2.0}{h} after the irradiation ends and
lasts \SI{10}{min}, with a detector that records $10\%$ of the decays. How many counts are recorded?
\challengeans{\emph{Production.} $\Act=\Prod(1-e^{-\lambda t})=\Prod(1-2^{-t/\Thalf})$; after two half-lives of
irradiation $\Act=10^{8}\times(1-\tfrac14)=\SI{7.5e7}{\becquerel}$. \emph{Waiting.} Two hours later (one
half-life) $\Act=\SI{3.75e7}{\becquerel}$. \emph{Counting.} The activity falls slightly during the count, so
integrate: decays $=\int_0^{600}\Act e^{-\lambda t}\,dt=(\Act/\lambda)(1-e^{-\lambda\times600\,\mathrm s})$, with
$\lambda=0.693/7200=\SI{9.63e-5}{s^{-1}}$; this gives $3.90\times10^{11}\times0.0561=2.19\times10^{10}$ decays.
Counts $=0.10\times2.19\times10^{10}=2.2\times10^{9}$. (Treating the activity as constant would give
$2.25\times10^{9}$, $3\%$ high.)}
\end{challengebox}

\begin{challengebox}{Why radium is warm (\chsref{ch:decay}{ch:decaylaw})}
Pierre Curie and Albert Laborde found in 1903 that radium salts are always warmer than their surroundings.
(a) Find the heat output of \SI{1.00}{g} of radium-226 alone ($Q=\SI{4.87}{\MeV}$, $\Act=\SI{3.66e10}{\becquerel}$)
in watts and in joules per hour. (b) A radium sample several weeks old gives five to six times this output.
Explain.
\challengeans{(a) $P=\Act Q=3.66\times10^{10}\times4.87\times1.602\times10^{-13}=\SI{0.029}{W}$, that is,
\SI{103}{J} per hour, without end on any human time scale ($\Thalf=\SI{1600}{\year}$). (b) Within a few weeks the
short-lived daughters (radon-222, polonium-218, lead-214, bismuth-214, polonium-214) reach secular
equilibrium with the radium. Each then has the same activity as the radium and adds its own decay energy:
three more $\alpha$ decays of $5.6$ to \SI{7.8}{\MeV} and two $\beta$ decays with their gamma rays, five to six
times the radium's own output in all (the neutrinos' share escapes). (Lead-210, with $\Thalf=\SI{22}{\year}$, holds the chain back from going
further on this time scale.)}
\end{challengebox}

\begin{challengebox}{Counting helium atoms (\chsref{ch:decay}{ch:decaylaw})}
Rutherford and Royds showed in 1909 that $\alpha$ particles are helium nuclei by collecting the gas they
became. What volume of helium, at standard temperature and pressure (\SI{22.4}{L} per mole, \chref{ch:kinetic}),
does \SI{1.00}{g} of radium-226 alone ($\Act=\SI{3.66e10}{\becquerel}$) produce in one year?
\challengeans{Each decay makes one helium atom. Decays in one year $=3.66\times10^{10}\times3.156\times10^{7}
=1.15\times10^{18}$ atoms, that is, $1.15\times10^{18}/6.022\times10^{23}=\SI{1.9e-6}{mol}$, or
$1.9\times10^{-6}\times\SI{22400}{cm^3}=\SI{0.043}{cm^3}$. Small but measurable, which is how the identification
was made. (The radium's activity hardly changes in a year, since $\Thalf=\SI{1600}{\year}$; its daughters in
equilibrium would add three more helium atoms per radium decay.)}
\end{challengebox}

\begin{challengebox}{The natural reactor at Oklo (\chsref{ch:decaylaw}{ch:reactions})}
Today natural uranium contains $0.720\%$ uranium-235 ($\Thalf=\SI{7.04e8}{\year}$) and $99.27\%$ uranium-238
($\Thalf=\SI{4.468e9}{\year}$). Find the percentage of uranium-235 in natural uranium 1.7 billion years ago.
At Oklo in Gabon, natural uranium deposits at that time sustained chain reactions with ordinary groundwater as
the moderator. Explain why that is impossible today.
\challengeans{Each isotope decays independently, so going back in time multiplies each by $e^{\lambda t}$, and the
ratio by $e^{(\lambda_{235}-\lambda_{238})t}$. With $\lambda_{235}=\SI{9.85e-10}{\year^{-1}}$ and
$\lambda_{238}=\SI{1.55e-10}{\year^{-1}}$, $(\lambda_{235}-\lambda_{238})t=8.29\times10^{-10}\times1.7\times10^{9}
=1.41$, so
\[
\frac{N_{235}}{N_{238}}=\frac{0.720}{99.27}\,e^{1.41}=0.00725\times4.10=0.0297,
\]
that is, $0.0297/1.0297=2.9\%$ uranium-235, similar to modern enriched reactor fuel. Today's $0.72\%$ is too
dilute: ordinary water absorbs enough neutrons that the multiplication factor $k$ cannot reach $1$ with natural
uranium, which is why light-water reactors need enriched fuel (\chref{ch:reactions}).}
\end{challengebox}

\section*{Looking Back on Part B}
\addcontentsline{toc}{section}{Looking Back on Part B}

\textbf{Part B turned the nucleus from an object of study into a source of energy and a tool.}
\begin{itemize}
\item \textbf{Nuclear reactions (\chref{ch:reactions}).} The bookkeeping of decay extends to reactions
X$(a,b)$Y: the $Q$-value is the rest energy lost, and an endoergic reaction needs a threshold kinetic energy
larger than $|Q|$, because momentum conservation forces some energy to stay as motion of the centre of mass
(CM), exactly the inelastic-collision argument of \chref{ch:collisions}. The binding-energy curve of
\chref{ch:nucleus}, peaked near iron, then explains both great energy sources at once: a heavy nucleus releases
about \SI{200}{\MeV} by splitting, light nuclei release more per nucleon by fusing. Fission becomes a power
source through the chain reaction, held exactly critical ($k=1$) with moderators that slow neutrons by
collisions and delayed neutrons that make control possible. Fusion is blocked by the Coulomb barrier and
proceeds in the Sun only because thermal motion (\chref{ch:kinetic}) and quantum tunnelling together let a tiny
fraction of nuclei through.
\item \textbf{Applications (\chref{ch:nucapps}).} The decay law read backwards dates bones with carbon-14 and the
solar system with uranium; inside the body it combines with excretion to give an effective half-life; in matter
it reappears as exponential attenuation, with the half-value layer playing the half-life. Dose is energy per
kilogram, weighted for how densely each radiation ionises, and protection rests on time, distance and
shielding. A simulation with nothing but random numbers reproduces the law and shows it to be a law of large
numbers.
\end{itemize}

\section*{Looking Back on the Volume}
\addcontentsline{toc}{section}{Looking Back on the Volume}

\textbf{Part A built the nucleus and the law of its decay.}
\begin{itemize}
\item \textbf{The nucleus (\chref{ch:nucleus}).} A nucleus of $Z$ protons and $A-Z$ neutrons is a tiny drop of
nuclear matter, of radius $R_0A^{1/3}$ and the same enormous density for every nucleus. It weighs less than its
separated parts, and the deficit times $c^2$ is the binding energy. Binding energy per nucleon rises steeply,
peaks near iron and then falls slowly: one curve that already contains fission, fusion and the limits of
stability.
\item \textbf{Radioactive decay (\chref{ch:decay}).} Unstable nuclei move towards the valley of stability by
$\alpha$ decay, $\beta^\pm$ decay and electron capture, often followed by $\gamma$ emission. Each decay obeys the
conservation of energy and momentum from Volume~I, together with conservation of charge, nucleon number and
lepton number, and its $Q$-value comes from atomic masses. Momentum conservation shares the energy: sharply between two
bodies in $\alpha$ decay, continuously among three in $\beta$ decay, which is how the neutrino was inferred
before it was seen.
\item \textbf{The decay law (\chref{ch:decaylaw}).} Each nucleus has a fixed probability per unit time of
decaying, independent of its age. From that one assumption follow $N=N_0e^{-\lambda t}$, the half-life and mean
life, activity, branching, and the chains that settle into secular or transient equilibrium. It is the same
first-order equation as chemical kinetics, and the same law that governs a hot body cooling by Newton's law
of cooling (\chref{ch:thermal}).
\end{itemize}

\textbf{Part B put the nucleus to work}: reactions, fission and fusion in \chref{ch:reactions}; dating,
medicine, dose and simulation in this chapter.

\medskip
\textbf{The volume in five sentences.} The nucleus is a tiny drop of uniform nuclear matter whose mass is
less than that of its parts, the deficit being its binding energy. The binding energy per nucleon peaks near
iron, so heavy nuclei can release energy by splitting and light ones by fusing. Unstable nuclei move towards the
valley of stability by $\alpha$, $\beta^\pm$ and electron-capture decays, often followed by $\gamma$ emission,
with $Q$-values fixed by atomic masses. Each nucleus decays at random with a fixed probability per unit time,
which makes populations decay exponentially and lets chains settle into equilibrium. That one exponential law
dates rocks and bones, schedules medical isotopes, shields workers and emerges unaided from a few lines of
random-number code.

\medskip
\textbf{The volume and the book.} Nothing in this volume needed a new principle of mechanics. Energy
conservation (\chref{ch:energy}), extended by $E=mc^2$ to include rest energy, gives every $Q$-value; momentum
conservation (\chref{ch:collisions}) shares that energy between products, fixes the threshold of a reaction and
explains the slowing of neutrons in a moderator; the Coulomb potential energy of two nuclei has the same
$1/r$ form as the gravitational potential energy of \chref{ch:gravitation}, with repulsion in place of
attraction. From Volume~II came the statistical view: the temperatures and speed distributions of
\chref{ch:kinetic} decide how often nuclei in the Sun meet fast enough to fuse, and the mean free path of a
gas molecule reappears as the mean distance $1/\mu$ a gamma ray travels in lead. What the nucleus added is
the idea that a law of nature can be exact about populations while saying nothing about individuals. The
exponential that began, in \chref{ch:toolkit}, as a mathematical function has become, by the end of the book,
the signature of pure chance.
